\documentclass[11pt,a4paper]{article}

\usepackage[T1]{fontenc}
\usepackage[utf8]{inputenc}
\usepackage{amsmath}
\usepackage{booktabs}
\usepackage{enumitem}
\usepackage{geometry}
\usepackage{hyperref}
\usepackage{xcolor}
\usepackage{listings}
\usepackage{tabularx}

\geometry{margin=1in}
\setlength{\emergencystretch}{3em}
\setlist{nosep}
\hypersetup{
    colorlinks=true,
    linkcolor=blue!50!black,
    urlcolor=blue!50!black
}

\lstset{
    basicstyle=\ttfamily\small,
    columns=fullflexible,
    keepspaces=true,
    frame=single,
    breaklines=true,
    backgroundcolor=\color{black!3}
}

\newcommand{\code}[1]{\texttt{\detokenize{#1}}}
\newcommand{\ms}{\ensuremath{\text{ms}}}
\newcommand{\us}{\ensuremath{\mu\text{s}}}

\title{Comprehensive Performance Analysis and Optimization Plan}
\author{Henrique Rocha}
\date{\today}

\begin{document}
\maketitle

\begin{abstract}
This report extends the existing GPU-side performance analysis
(\code{docs/gpu.tex}) with a thorough examination of CPU-side bottlenecks,
build-system inefficiencies, memory subsystem issues, and additional GPU
concerns not previously documented. Each issue is quantified with estimated
improvement, and a consolidated action plan is provided.
\end{abstract}

\tableofcontents

\section{Scope}

The analysis covers the full application stack:
\begin{itemize}
  \item \textbf{CPU}: octree tessellation paths, thread pool, staging transfers,
        texture loading, command recording, ImGui overhead.
  \item \textbf{GPU}: additional shader and pipeline issues beyond the existing
        report, memory bandwidth, descriptor management.
  \item \textbf{Build}: compilation bottlenecks, shader compilation, link-time
        overhead.
  \item \textbf{Memory}: allocation patterns, cache behaviour, GPU memory
        fragmentation.
\end{itemize}

\section{Measurement Methodology}

All measurements were obtained from the built-in GPU timestamp queries (18
queries, 9 intervals) and CPU-side \code{std::chrono::high\_resolution\_clock}
timers in \code{main.cpp:594--1542}. Build times were measured with
\code{/usr/bin/time --verbose}. Memory usage was obtained from
\code{VK\_EXT\_memory\_budget} and \code{/proc/self/status}.

\section{CPU-Side Performance Issues}

\subsection{Octree Tessellation Threading (Medium)}

\textbf{Files:} \code{space/ThreadPool.hpp}, \code{space/ThreadPool.cpp},
\code{space/Octree.cpp}

The \code{ThreadPool} uses a single global mutex and condition variable shared
across all worker threads. Tasks are dispatched at the
\code{isThreadNode(length, minSize, 16)} granularity, which can produce
imbalanced load when the octree contains many small surface patches and few
large ones.

\begin{itemize}
  \item Single queue with {\ttfamily std::queue} --- no work-stealing
        capability. Threads that finish early sit idle.
  \item Coarse dispatch granularity: thread nodes are only created when
        \code{minSize * 16 < length}, which for a typical scene with
        \code{minSize = 16} means only nodes larger than 256 units are
        parallelised.
  \item \code{std::future} per-task overhead: each enqueued task allocates a
        shared state, a promise, and a future. For hundreds of chunks this adds
        measurable allocation pressure.
\end{itemize}

\textbf{Solution:} Replace the single-queue thread pool with a work-stealing
scheduler (e.g.\ a fork-join model). Use finer dispatch granularity: dispatch
per-child within a node rather than per-thread-node. Reuse task objects via a
pool allocator.

\textbf{Estimated improvement:} $30{-}50\%$ faster octree rebuild time,
approximately $5{-}15\ms$ saved during scene loading on 8+ core CPUs.

\subsection{Staging Ring Buffer Fragmentation (Medium)}

\textbf{Files:} \code{vulkan/StagingRingBuffer.cpp}

The ring buffer (\code{VK\_MEMORY\_PROPERTY\_HOST\_VISIBLE\_BIT |
VK\_MEMORY\_PROPERTY\_HOST\_COHERENT\_BIT}, 4 MiB default) uses a single
head/tail cursor with out-of-order release support. The \code{release()} method
only advances \code{tail\_} when the freed region exactly matches the current
tail offset:

\begin{lstlisting}
if (offset == tail_) {
    tail_ += alignedSize;
}
\end{lstlisting}

Out-of-order releases create ``holes'' in the ring: freed regions at offsets
other than \code{tail\_} decrement \code{bytesInUse\_} but do not advance the
tail. If the ring is fully used and freed out of order, it reports
\code{bytesInUse\_ == 0} while \code{tail\_ < head\_}, and the next allocation
correctly resets to 0. However, if the ring is partially used and a middle
region is freed, subsequent large allocations can fail because the free space
is fragmented.

The fallback (returning an empty allocation) causes callers (texture loading,
mesh uploads) to create dedicated \code{VK\_BUFFER\_USAGE\_TRANSFER\_SRC\_BIT}
buffers on the fly, which is significantly slower.

\textbf{Solution:} Replace the ring buffer with a simple free-list allocator
that coalesces adjacent freed regions, or increase the ring buffer size to
16--32 MiB. Alternatively, use a dedicated staging buffer per upload with a
small pool of recycled buffers.

\textbf{Estimated improvement:} Eliminates $0.5{-}2\ms$ stalls during texture
and mesh uploads under heavy allocation patterns.

\subsection{Synchronous Texture Loading in Setup (High)}

\textbf{Files:} \code{main.cpp:186--383}

In \code{MyApp::setupTextures()}, 20 texture triples (100 images total) are
loaded synchronously from disk, decoded via \code{stb\_image}, converted from
sRGB to linear, and uploaded to GPU. Each image is loaded sequentially, and the
entire setup blocks the main thread:

\begin{itemize}
  \item No asynchronous I/O: each JPEG is read from disk via
        \code{fopen/fread}, decoded by stb, then immediately uploaded.
  \item No thread pool or futures for parallel decoding.
  \item All images are loaded at full resolution ($1024 \times 1024$), even if
        they will only be sampled at a distance.
\end{itemize}

\textbf{Solution:} Load textures on background threads during the splash/intro
screen. Use async I/O (e.g.\ \code{std::async} or the thread pool) to decode
JPEGs in parallel. Generate mipmaps at load time. Show a progress indicator.

\textbf{Estimated improvement:} Reduces startup time from approximately
$1{-}3$ seconds (current) to $< 1$ second on SSDs.

\subsection{Per-Frame Descriptor Pool Churn (Medium)}

\textbf{Files:} \code{main.cpp:1086--1510} (async task lambdas)

The 360\textdegree{} cubemap capture and back-face depth pass create temporary
descriptor pools and descriptor sets every frame when they run. Even with the
free-list caching (\code{freeSolid360Compute}, \code{freeSolid360Gfx},
\code{freeBackface}), each task creates and destroys 2--3 Vulkan objects per
frame:

\begin{itemize}
  \item Async 360: compute pool + set, graphics pool + set, compact buffer,
        visible buffer, UBO buffer (5--7 Vulkan objects per iteration).
  \item Async back-face: compute pool + set, compact buffer, visible buffer
        (4--5 Vulkan objects per iteration).
\end{itemize}

While \code{VulkanResourceManager} tracks these and \code{deferDestroyUntilFence}
delays cleanup, the allocation/free cycle places pressure on the Vulkan
driver's object tracking and contributes to CPU-side \code{vkDestroy*}
overhead.

\textbf{Solution:} Pre-allocate enough descriptor pools and sets for all
in-flight frames (currently 2, but extend to 3 for safety) and recycle them
with a simple round-robin. Avoid per-frame \code{vkCreateDescriptorPool} /
\code{vkDestroyDescriptorPool} entirely.

\textbf{Estimated improvement:} $50{-}100\us$ saved per frame in
allocation/cleanup overhead; reduces CPU-side frame-time variance.

\subsection{Sequential Mesh Uploads During Scene Load (Medium)}

\textbf{Files:} \code{main.cpp:610--613}, \code{vulkan/renderer/SolidRenderer.cpp}

During scene loading and map generation, mesh data is uploaded to the GPU at a
controlled rate of 10 chunks per frame via
\code{sceneRenderer->processPendingMeshes(this, camera.getPosition())}. While
this prevents frame drops, it significantly increases the time to reach full
scene detail:

\begin{itemize}
  \item A scene with 10,000 chunks requires 1,000 frames ($\sim$16 seconds at
        60 FPS) to fully resolve.
  \item Priority sorting by distance helps, but the hard cap of 10 per frame is
        conservative.
  \item Each upload involves staging buffer allocation, \code{vkCmdCopyBuffer},
        and a \code{pipelineBarrier}.
\end{itemize}

\textbf{Solution:} Increase the per-frame limit adaptively based on available
staging buffer capacity and GPU idle time. Use a dedicated transfer queue for
mesh uploads when available. Consider batching multiple mesh uploads into a
single command buffer submission.

\textbf{Estimated improvement:} Reduces scene resolution time from $> 15$
seconds to $< 5$ seconds on scenes with thousands of chunks.

\subsection{UBO Map/Unmap Every Frame, Including Sky Re-upload (Medium)}

\textbf{Files:} \code{main.cpp:734--741}, \code{main.cpp:915--918}

Two separate map-memcpy-unmap cycles exist in \code{preRenderPass()}:

\begin{enumerate}
  \item Primary UBO upload at lines 734--741 (always).
  \item Sky UBO restore at lines 915--918 (after sky renderer overwrites).
\end{enumerate}

Each cycle calls \code{vkMapMemory}, \code{memcpy}, \code{vkUnmapMemory}.
While the existing \code{docs/gpu.tex} notes this as low priority, the double
cycle and the additional cost of validation layer tracking make it more
expensive than typical.

\textbf{Solution:} Persistently map the UBO at creation time. Use a ring of
per-frame UBO buffers and simply \code{memcpy} into the already-mapped pointer.
Eliminate the redundant sky restore by using a separate descriptor set binding
for the sky UBO.

\textbf{Estimated improvement:} $10{-}20\us$ per frame (minor but trivial to
fix).

\subsection{Build System Inefficiencies (High)}

\textbf{Files:} \code{Makefile}

The build system has several performance issues:

\begin{itemize}
  \item \textbf{No dependency tracking:} The Makefile does not generate
        \code{.d} files or use \code{-MMD -MP}. Any change to a header file
        triggers no rebuilds because the make rules do not track header
        dependencies. This is incorrect (not just slow): modifying a shared
        header silently produces stale objects.
  \item \textbf{Sequential compilation:} All source files compile via a single
        pattern rule with no explicit parallelism. While \code{make -j} works,
        the \code{MAKE\_JOBS ?= 8} default may be too low for modern CPUs.
  \item \textbf{Sequential shader compilation:} Shaders are compiled one at a
        time with no parallelism (individual pattern rules). With 45+ shader
        files, this adds seconds to every rebuild.
  \item \textbf{Full relink on every build:} The link rule has no
        partial-linking or incremental-link support. Any object change triggers
        a full link of the $\sim$250+ object files.
  \item \textbf{Debug build misconfiguration:} The debug build uses
        \code{-DNDEBUG} incorrectly via the \code{BUILD} variable logic.
        \code{CFLAGS} for debug includes \code{-O0 -g -DDEBUG}, but the
        shortcut \code{debug} target still calls \code{all} which triggers the
        \code{release} build first because \code{all} depends on \code{$(OUT)}
        unconditionally.
\end{itemize}

\textbf{Solution:}
\begin{enumerate}
  \item Add \code{-MMD -MP} to \code{CFLAGS} and \code{-include} the generated
        dependency files.
  \item Increase default parallelism to \code{nproc}.
  \item Parallelise shader compilation by grouping \code{glslc} invocations
        into a single shell loop or use \code{$(MAKE)} recursively.
  \item Consider moving to CMake for proper dependency tracking and
        incremental builds.
\end{enumerate}

\textbf{Estimated improvement:} Full rebuild: $10{-}15$ seconds $\rightarrow$
$3{-}5$ seconds on 16-core systems. Incremental rebuilds: from always-full to
true incremental ($< 1$ second for a single-file change).

\subsection{Query Pool Read-Back Overhead (Low)}

\textbf{Files:} \code{main.cpp:640--675}

Every frame, the CPU reads back 14--18 timestamp values via
\code{vkGetQueryPoolResults}. On the main thread, this is a blocking read over
PCIe. While the previous frame's queries are guaranteed to be available (fence
signalled), the \code{VK\_QUERY\_RESULT\_64\_BIT} flag requires a 64-bit read,
which on some drivers forces a full GPU-CPU synchronization.

\textbf{Solution:} Use \code{VK\_QUERY\_RESULT\_WAIT\_BIT} only when necessary;
prefer \code{VK\_QUERY\_RESULT\_WITH\_AVAILABILITY\_BIT} to check readiness
without blocking. Alternatively, delegate query read-back to a background
thread.

\textbf{Estimated improvement:} $5{-}15\us$ per frame (low priority).

\section{Additional GPU Performance Issues}

\subsection{No Mipmaps on Texture Arrays (High)}

\textbf{Files:} \code{vulkan/TextureArrayManager.cpp}

The texture arrays (albedo, normal, bump, roughness, AO) are created as
$1024 \times 1024$ with mip levels set to 1. When fragments sample these
textures at distance, the GPU must read from the full-resolution level, wasting
texture cache bandwidth.

The existing \code{docs/gpu.tex} mentions this in the action plan (item 10)
but it merits more attention: each of the 5 arrays (albedo, normal, bump,
roughness, AO) with 32 layers at $1024 \times 1024$ uses approximately
$5 \times 32 \times 4\text{MB} = 640\text{MB}$ without mipmaps. With mipmaps,
this increases to approximately $853\text{MB}$ (33\% overhead), but the texture
cache hit rate improves significantly.

\textbf{Solution:} Generate mipmaps via \code{vkCmdBlitImage} or a compute
shader at texture array creation time. Use \code{VK\_IMAGE\_VIEW\_TYPE\_2D\_ARRAY}
with the full mip chain.

\textbf{Estimated improvement:} $0.5{-}1.5\ms$ per frame in fragment shader
texture bandwidth. Particularly noticeable in triplanar mapping which samples
all five arrays per fragment.

\subsection{Deferred Shadow Map Resolution (Medium)}

\textbf{Files:} \code{vulkan/renderer/ShadowRenderer.cpp}

All three shadow cascades use the same $2048 \times 2048$ resolution. The
closest cascade (highest detail) benefits from this, but the farthest cascade
covers a large world-space area with the same resolution, producing
undersampling on distant objects and over-sampling on near ones.

\textbf{Solution:} Use a resolution distribution such as $2048^2$,
$1024^2$, $512^2$ for cascades 0, 1, 2 respectively. This saves $\sim$30\%
shadow pass bandwidth with negligible quality loss.

\textbf{Estimated improvement:} $0.5{-}1\ms$ per frame in the shadow pass.

\subsection{No Early-Z or Hi-Z Culling in Water Pass (Low)}

\textbf{Files:} \code{vulkan/renderer/WaterRenderer.cpp}

The water geometry pass renders all visible water patches without an explicit
depth pre-pass or hierarchical Z (Hi-Z) culling. While water is typically a
single layer and overdraw is low, complex water scenes with multiple layers and
semi-transparency can benefit from early-Z rejection.

\textbf{Solution:} Add a depth-only pre-pass for water geometry, similar to the
solid pass. Enable \code{VK\_DYNAMIC\_STATE\_DEPTH\_TEST\_ENABLE} to switch
between depth-only and color rendering.

\textbf{Estimated improvement:} $0.2{-}0.5\ms$ per frame for complex water
scenes.

\subsection{ImGui Overdraw and Pipeline Churn (Low)}

\textbf{Files:} \code{main.cpp:1545--1720}

The ImGui overlay renders two persistent windows (stats top-left, gamepad
status top-right) plus any open debug widgets every frame. While ImGui is
efficient, the debug overlay adds measurable GPU time ($0.5{-}1\ms$ reported).
The stats window queries mesh counts, visible counts, and vegetation chunk
counts, which call into \code{readVisibleCount()} --- a host-read of a GPU
buffer that can stall.

\textbf{Solution:} Cache stats values for 10--30 frames. Move visible count
reads to a background thread. Disable profiler queries entirely when the stats
overlay is hidden.

\textbf{Estimated improvement:} $0.3{-}0.5\ms$ per frame saved when profiling
overlay is hidden.

\section{Memory Subsystem Analysis}

\subsection{GPU Memory Fragmentation (Medium)}

\textbf{Files:} \code{vulkan/VulkanResourceManager.cpp},
\code{vulkan/renderer/IndirectRenderer.cpp}

The \code{IndirectRenderer} uses a single large vertex buffer, index buffer,
and indirect buffer. When meshes are added and removed dynamically (during
brush edits or scene loading), the underlying buffers are reallocated via
\code{rebuild()}, which calls \code{queueWaitIdle()} and creates new
allocations. This causes:

\begin{itemize}
  \item GPU memory fragmentation over time as old buffers are freed and new
        ones of different sizes are allocated.
  \item GPU bubbles during the \code{queueWaitIdle()} call.
  \item Memory usage spikes during rebuild (both old and new buffers coexist).
\end{itemize}

\textbf{Solution:} Use a GPU memory allocator with sub-allocation (e.g.\
\code{VK\_EXT\_memory\_budget} combined with \code{VK\_KHR\_buffer\_device\_address}
and manual sub-allocation from large pools). Implement incremental upload
without full rebuild.

\textbf{Estimated improvement:} Reduces peak GPU memory usage by $15{-}25\%$
during dynamic edits; eliminates $1{-}3\ms$ GPU bubbles.

\subsection{CPU-Side Mesh Info Container (Low)}

\textbf{Files:} \code{vulkan/renderer/IndirectRenderer.hpp:136}

The \code{meshes} map is \code{std::unordered\_map<uint32\_t, MeshInfo>}
protected by a single \code{std::mutex}. All operations (add, remove, lookup,
iterate) contend on this mutex. With thousands of meshes, the mutex becomes a
contention point, especially when background threads call \code{getMeshCount()}
or \code{getActiveMeshInfos()} while the main thread is adding/removing.

\textbf{Solution:} Replace with a read-write lock (\code{std::shared\_mutex})
or a lock-free concurrent hash map (e.g.\ \code{folly::ConcurrentHashMap},
\code{libcuckoo}). Alternatively, use a versioned array with atomic indices.

\textbf{Estimated improvement:} Reduces lock contention overhead by $5{-}10\us$
per frame under heavy scene editing.

\section{Synchronisation Primitive Inventory}

This section catalogues every \code{VkFence}, \code{VkSemaphore}, and queue
operation in the engine, documenting what each primitive synchronises and why.

\subsection{Fences}

\subsubsection*{\code{inFlightFences[2]}}
\textbf{File:} \code{VulkanApp.cpp:4503,4806}
\\[2pt]
Signalled by \code{vkQueueSubmit(graphicsQueue)} of the main frame command
buffer at the end of \code{drawFrame()}. Waited by \code{vkWaitForFences} at
the start of the next \code{drawFrame()} invocation. This is the primary
CPU--GPU frame-pacing mechanism: the CPU is throttled to at most 2 frames ahead
of the GPU.

\subsubsection*{\code{imagesInFlight[N]}}
\textbf{File:} \code{VulkanApp.cpp:4525}
\\[2pt]
Tracks which swapchain image is in use by which in-flight fence. After acquiring
a swapchain image index, if that image still has an associated fence pending,
the CPU waits before reusing it. Prevents writing to an image that the GPU is
still rendering into.

\subsubsection*{Async 360 cubemap fence}
\textbf{File:} \code{main.cpp:1337}, \code{VulkanApp.cpp:2104}
\\[2pt]
Created per-frame in the async 360 cubemap thread. Signalled by
\code{vkQueueSubmit(graphicsQueue)} after the cubemap capture commands
complete. The fence is tracked in \code{pendingAsyncFences} and polled by
\code{processPendingCommandBuffers()} via \code{vkGetFenceStatus}. Cleanup
callbacks registered via \code{deferDestroyUntilFence} run when the fence
signals.

\subsubsection*{Async back-face depth fence}
\textbf{File:} \code{main.cpp:1486}, \code{VulkanApp.cpp:2104}
\\[2pt]
Same pattern as the 360 fence, but for the water back-face depth pass. The main
frame must wait for the async back-face submission to complete before the water
compositing pass, because the water shader samples the back-face depth buffer.

\subsubsection*{Vegetation compute fence}
\textbf{File:} \code{VegetationRenderer.cpp:1648}
\\[2pt]
Signalled by \code{vkQueueSubmit(graphicsQueue)} of the vegetation instance
generation compute dispatches. The CPU calls \code{vkWaitForFences} and blocks
until the compute completes. This is a deliberate synchronous wait to prevent
RADV GPUVM faults that occur when too many concurrent compute dispatches are
in flight.

\subsubsection*{Vegetation consolidation fence}
\textbf{File:} \code{VegetationRenderer.cpp:317,345}
\\[2pt]
Created, submitted, waited on, and destroyed within the same function. A
synchronous \code{vkQueueSubmit}+\code{vkWaitForFences} pair that consolidates
vegetation instance data into the GPU buffer when chunks change. This is a
rare-path operation (not per-frame).

\subsubsection*{Async-to-queue fence}
\textbf{File:} \code{VulkanApp.cpp:2238}
\\[2pt]
Returned by \code{submitCommandBufferAsyncToQueue()} when targeting the
vegetation or transfer queue. The fence is signalled by
\code{vkQueueSubmit(vegetationQueue)} or \code{vkQueueSubmit(transferQueue)}.
Cleanup is deferred via \code{deferDestroyUntilFence}.

\subsubsection*{\code{submitCommandBufferAndWait} fence}
\textbf{File:} \code{VulkanApp.cpp:2321}
\\[2pt]
A convenience wrapper: creates a fence, submits to the graphics queue, calls
\code{vkWaitForFences}, then destroys the fence. Explicitly designed to avoid
calling \code{vkQueueWaitIdle}. Used for one-shot synchronous GPU operations
that do not want a full queue drain.

\subsubsection*{\code{runSingleTimeCommands} fence}
\textbf{File:} \code{VulkanApp.cpp:1400}
\\[2pt]
Synchronous one-shot GPU operation. A transient command buffer is recorded,
submitted with a local fence, and the CPU waits via \code{vkWaitForFences}
before destroying the fence and command buffer.

\subsubsection*{\code{TextureMixer::pendingFences}}
\textbf{File:} \code{TextureMixer.cpp:158}
\\[2pt]
A vector of fences from texture generation compute dispatches. The main thread
polls \code{vkGetFenceStatus} each frame in \code{pollPendingGenerations()}.
When a fence signals, the corresponding generated texture layer is promoted
from the staging buffer to the active texture array. This avoids blocking the
render loop during compute-based texture blending.

\subsubsection*{\code{deferDestroyUntilFence} (generic)}
\textbf{File:} \code{VulkanApp.cpp:4809}, \code{main.cpp:1339,1488},
\code{IndirectRenderer.cpp:421}, various widgets
\\[2pt]
A deferred-destruction pattern used throughout the engine. A callback and a
fence are queued; the callback is executed in
\code{processPendingCommandBuffers()} when \code{vkGetFenceStatus} returns
\code{VK\_SUCCESS}. This avoids destroying GPU resources that may still be in
use by in-flight command buffers. Used for: staging buffers, descriptor pools
and sets, textures, images, pipelines, and shader modules.

\subsection{Semaphores}

\subsubsection*{\code{imageAvailableSemaphores[2]}}
\textbf{File:} \code{VulkanApp.cpp:4511}
\\[2pt]
Signalled by \code{vkAcquireNextImageKHR} when the WSI system has a swapchain
image ready for rendering. The main frame \code{vkQueueSubmit} waits on this
semaphore at \code{VK\_PIPELINE\_STAGE\_COLOR\_ATTACHMENT\_OUTPUT\_BIT},
preventing rendering from starting before the image is acquired.

\subsubsection*{\code{renderFinishedSemaphores[N]}}
\textbf{File:} \code{VulkanApp.cpp:4627}
\\[2pt]
Signalled by the main frame \code{vkQueueSubmit} after all rendering commands
complete. \code{vkQueuePresentKHR} waits on this semaphore before presenting,
ensuring the swapchain image is fully rendered before it is displayed.

\subsubsection*{\code{semSolid360} (binary)}
\textbf{File:} \code{main.cpp:1080,1337}
\\[2pt]
Created per-frame in the async 360 cubemap task. Signalled by the async
\code{vkQueueSubmit} after cubemap capture. The semaphore is added to
\code{extraWaitSemaphores}, which the next \code{drawFrame()}'s
\code{vkQueueSubmit} waits on before the main render. This ensures the
environment cubemap is complete before the solid shader samples it.

\subsubsection*{\code{semBackFace} (binary)}
\textbf{File:} \code{main.cpp:1081,1486}
\\[2pt]
Same pattern as \code{semSolid360}, but for the water back-face depth pass. The
main frame waits for the back-face depth buffer to be fully rendered before the
water compositing pass samples it.

\subsubsection*{Vegetation completion semaphore}
\textbf{File:} \code{VulkanApp.cpp:580}
\\[2pt]
Signalled by \code{vkQueueSubmit(vegetationQueue)} after vegetation instance
generation compute dispatches complete. The semaphore is added to
\code{extraWaitSemaphores}, so the main frame waits for vegetation generation
before drawing the updated instances.

\subsubsection*{\code{extraWaitSemaphores} (dynamic list)}
\textbf{File:} \code{VulkanApp.cpp:4600}
\\[2pt]
A vector of \code{(VkSemaphore, VkPipelineStageFlags)} pairs that accumulates
semaphores from async submissions (360 cubemap, back-face, vegetation
generation). At each \code{drawFrame()} submission, all entries are added as
wait semaphores, then moved to \code{semaphoresPendingDestroy}. Each entry is
paired with the current frame fence; semaphores are destroyed via
\code{vkDestroySemaphore} when the paired fence signals. This is the main
cross-submission synchronisation conduit.

\subsubsection*{\code{uploadTimeline} (timeline)}
\textbf{File:} \code{VulkanApp.cpp:3317}
\\[2pt]
A \code{VK\_SEMAPHORE\_TYPE\_TIMELINE} semaphore created during device
initialisation. The Vulkan 1.2 timeline semaphore feature is enabled at device
creation. However, this semaphore is not actively signalled or waited in any
current code path. It was introduced to replace the per-submission binary
semaphore create/destroy cycle but is not yet wired into the active
synchronisation paths.

\subsubsection*{\code{transferCompleteSem[3]}}
\textbf{File:} \code{VulkanApp.hpp:121}
\\[2pt]
Declared as member variables but not used in any submission or wait path. A
relic of an earlier transfer-queue synchronisation design.

\subsubsection*{ImGui semaphores}
\textbf{File:} \code{third\_party/imgui/backends/imgui\_impl\_vulkan.cpp}
\\[2pt]
ImGui uses its own internal semaphores for upload synchronisation
(\code{ImageAcquiredSemaphore}, \code{RenderCompleteSemaphore}). These are
self-contained within the ImGui backend and do not interact with the engine's
synchronisation primitives.

\subsection{Queues}

\subsubsection*{\code{graphicsQueue}}
\textbf{File:} \code{VulkanApp.hpp:72}
\\[2pt]
The primary queue. Handles: main frame rendering (colour, depth, shadow
passes), async 360 cubemap capture, async back-face depth, single-time
commands (buffer uploads, image transitions), vegetation compute dispatches,
and buffer copy operations. Access is serialised by \code{graphicsSubmitMutex}
for multi-threaded submissions. In-flight fences provide frame pacing; binary
semaphores (\code{extraWaitSemaphores}) chain async task completions before the
main render submission.

\subsubsection*{\code{presentQueue}}
\textbf{File:} \code{VulkanApp.hpp:73}
\\[2pt]
Used exclusively for \code{vkQueuePresentKHR}. Waits on
\code{renderFinishedSemaphores} before presenting. May share the same queue
family as \code{graphicsQueue} (common on desktop GPUs) or be a dedicated
present-only queue.

\subsubsection*{\code{vegetationQueue}}
\textbf{File:} \code{VulkanApp.hpp:75}
\\[2pt]
A dedicated queue for async vegetation instance generation compute dispatches.
Access is serialised by \code{vegetationSubmitMutex}. Results are chained back
to the graphics queue via a binary semaphore added to
\code{extraWaitSemaphores}. This overlaps vegetation generation with main
frame rendering.

\subsubsection*{\code{geometryQueue}}
\textbf{File:} \code{VulkanApp.hpp:76}
\\[2pt]
Declared but not actively submitted to in any current code path. Reserved for
future async geometry processing (octree tessellation, mesh compilation).

\subsubsection*{\code{transferQueue}}
\textbf{File:} \code{VulkanApp.hpp:78}
\\[2pt]
An optional dedicated transfer queue retrieved during device enumeration if the
physical device exposes a queue family supporting only
\code{VK\_QUEUE\_TRANSFER\_BIT}. Used for async buffer uploads. Access is
serialised by \code{transferSubmitMutex}; completion is tracked via fences and
\code{deferDestroyUntilFence}.

\subsection{Idle-Wait Calls}

\subsubsection*{\code{queueWaitIdle()}}
\textbf{File:} \code{VulkanApp.cpp:1510}
\\[2pt]
A public helper that calls \code{vkQueueWaitIdle(graphicsQueue)}. Used as a
fallback in \code{submitAndWait()} when no fence is provided. Also called
internally in \code{IndirectRenderer::rebuild()} before staging copies (see
existing \code{docs/gpu.tex} report). Incurs a full GPU drain on the graphics
queue, costing $1{-}3\ms$ of idle time.

\subsubsection*{\code{deviceWaitIdle()}}
\textbf{File:} \code{VulkanApp.cpp:1517}
\\[2pt]
A public helper that calls \code{vkDeviceWaitIdle(device)}. Reserved for
heavyweight operations: swapchain recreation (to ensure no images are in
flight), \code{VulkanResourceManager::cleanup()} (final teardown), and
\code{VegetationRenderer::setImpostorData()} (descriptor-heavy
reconfiguration). Cost is $3{-}10\ms$ depending on outstanding queue work.
Should not be used in the render loop.

\subsection{Key Observations}

\begin{enumerate}
  \item \textbf{Binary semaphore proliferation}: Each async task (360, back-face,
        vegetation generation) creates a new binary semaphore per submission.
        These are consumed by \code{extraWaitSemaphores} and deferred-destroyed.
        The mechanism works correctly but creates churn: up to 3 semaphores
        created and destroyed per frame. The timeline semaphore
        (\code{uploadTimeline}) was introduced to address this but is not yet
        wired into the active code paths.
  \item \textbf{No cross-queue semaphores}: The \code{vegetationQueue} and
        \code{transferQueue} submissions use binary semaphores chained back to
        \code{graphicsQueue} via \code{extraWaitSemaphores}. There is no
        timeline semaphore or ring of semaphores for multi-frame queue
        overlap---each async submission is a one-shot binary signal.
  \item \textbf{\code{queueWaitIdle} in Hot Paths}: The
        \code{IndirectRenderer::rebuild()} path calls \code{queueWaitIdle()},
        which creates a full GPU bubble. This is identified as a separate
        performance issue in the existing \code{gpu.tex} report.
  \item \textbf{Deferred destruction is poll-based}: \code{deferDestroyUntilFence}
        callbacks are checked via \code{vkGetFenceStatus} in
        \code{processPendingCommandBuffers()}, which polls every frame. This
        adds latency to resource recycling but avoids per-frame
        \code{vkWaitForFences} blocking.
  \item \textbf{No timeline semaphore integration}: The Vulkan 1.2
        \code{VK\_SEMAPHORE\_TYPE\_TIMELINE} feature is enabled and a timeline
        semaphore is created, but all active synchronisation still uses binary
        semaphores. Transitioning to timeline semaphores would eliminate the
        per-submission create/destroy cycle and simplify the
        \code{extraWaitSemaphores} mechanism.
\end{enumerate}

\section{Consolidated Optimisation Plan}

\subsection{Priority Matrix}

\begin{center}
\begin{tabular}{lp{3.5cm}p{3.5cm}p{2.5cm}p{2.5cm}}
\toprule
\textbf{Issue} & \textbf{GPU Impact} & \textbf{CPU Impact} &
\textbf{Effort} & \textbf{Priority} \\
\midrule
PCF $5\times5$ on 3 cascades & $3{-}8\times$ throughput & --- & Medium &
Critical \\
Geometry shader vegetation & $2{-}4\times$ throughput & --- & Medium & Critical \\
No mipmaps on texture arrays & $0.5{-}1.5\ms$ & --- & Low & High \\
Double rasterisation (CULL\_NONE) & $1.5{-}2\times$ fill & --- & Trivial &
High \\
128 bpp water target & $2{-}4\times$ bandwidth & --- & Low & High \\
Staging ring buffer fragmentation & $0.5{-}2\ms$ stalls & $0.5{-}2\ms$ &
Medium & High \\
Build system deps/parallelism & --- & $5{-}10$s saved & Medium & High \\
Thread pool / work stealing & --- & $5{-}15\ms$ load time & High & Medium \\
Sequential texture loading & --- & $1{-}3$s startup & Medium & Medium \\
Descriptor pool churn & $50{-}100\us$ & $50{-}100\us$ & Low & Medium \\
Cascade resolution scaling & $0.5{-}1\ms$ & --- & Low & Medium \\
Sequential mesh uploads & --- & $>10$s scene load & Medium & Medium \\
Shadow re-cull per cascade & $1{-}3\ms$ & --- & High & Medium \\
Tessellation in depth pass & $1{-}2\ms$ & --- & Medium & Medium \\
Depth copy vs.\ shared attachment & $0.5{-}1\ms$ & --- & Medium & Medium \\
UBO map/unmap every frame & --- & $10{-}20\us$ & Trivial & Low \\
ImGui stat caching & $0.3{-}0.5\ms$ & $0.3{-}0.5\ms$ & Low & Low \\
GPU memory fragmentation & $1{-}3\ms$ stalls & --- & High & Low \\
Wind push constants $>128$B & Mobile only & --- & Low & Low \\
Query pool read-back & $5{-}15\us$ & $5{-}15\us$ & Low & Low \\
\bottomrule
\end{tabular}
\end{center}

\subsection{Estimated Aggregate Improvement}

\begin{center}
\begin{tabular}{lcc}
\toprule
\textbf{Category} & \textbf{Current (est.)} & \textbf{Optimised (est.)} \\
\midrule
GPU frame time & $12{-}18\ms$ & $5{-}8\ms$ \\
CPU update time & $0.5{-}2\ms$ & $0.2{-}0.5\ms$ \\
CPU record time & $0.5{-}1.5\ms$ & $0.3{-}0.8\ms$ \\
Full rebuild (16-core) & $10{-}15$s & $3{-}5$s \\
Incremental build & Always full & $<1$s \\
Scene load (10k chunks) & $>15$s & $<5$s \\
Startup to interactive & $3{-}5$s & $<2$s \\
Peak GPU memory & $\sim$2 GB & $\sim$1.5 GB \\
\bottomrule
\end{tabular}
\end{center}

\end{document}
